\chapter{Welcome to CampusTrack}

\section{What is CampusTrack?}

CampusTrack is a complete software system for schools and colleges. It
connects four groups of people:

\begin{itemize}[itemsep=2pt]
  \item \textbf{Parents} -- see, on their own phone, when their child
        arrives at school and leaves school, which rooms and classes the
        child attended, the full semester timetable, and progress
        reports written by teachers. Parents can also send feedback to
        the school and receive replies.
  \item \textbf{Students} -- see their timetable, download assignments
        and notes by scanning a QR code, and upload their projects,
        activities and theses from their phone.
  \item \textbf{Teachers} -- use a website to write progress updates
        about students, publish assignments and notes, answer parent
        feedback, and review student uploads.
  \item \textbf{Administrators} -- manage everything: classes, sections,
        rooms, RFID readers, timetables and user accounts. Classes and
        sections are \emph{modular}: the admin can add or remove them at
        any time.
\end{itemize}

\section{How does the attendance work? (in plain words)}

Every student carries a normal-looking \textbf{ID card}. Inside the card
there is a tiny \textbf{RFID chip} -- like the chip in a contactless bus
card. The student does not need to touch anything or stand in a queue.

At the \textbf{school gate} and in \textbf{every room} (classrooms,
laboratories, library, discussion rooms, auditorium) there are small
\textbf{antennas} fixed above the door. When a student simply walks
through, the antennas notice the card -- automatically, from a distance.

\subsection{How the system knows IN from OUT}

The trick is the \emph{order} in which the antennas see the card:

\begin{itemize}
  \item Walking \textbf{in}, the card passes antenna 1 first, then 2,
        then 3 (at the gate).
  \item Walking \textbf{out}, the same antennas see the card in the
        reverse order: 3, then 2, then 1.
\end{itemize}

\begin{figure}[H]
\centering
\begin{tikzpicture}[font=\sffamily, >={Stealth[length=3mm]}]
  % gate structure
  \fill[black!10] (0,0) rectangle (1.0,4.2);
  \fill[black!10] (9.0,0) rectangle (10.0,4.2);
  \node at (5,4.6) {\textbf{School gate (3 antennas)}};
  % antennas
  \foreach \x/\n in {2.5/1, 5.0/2, 7.5/3}{
     \draw[fill=campusblue!20, draw=campusblue, thick] (\x-0.35,3.4) rectangle (\x+0.35,3.9);
     \node[text=campusblue] at (\x,3.65) {\small\textbf{A\n}};
     \draw[campusblue!60, dashed] (\x,3.4) -- (\x-0.7,1.2);
     \draw[campusblue!60, dashed] (\x,3.4) -- (\x+0.7,1.2);
  }
  % entry arrow
  \draw[->, very thick, okgreen] (1.3,2.1) -- (8.7,2.1)
     node[midway, above=2pt, text=okgreen] {\textbf{ENTRY}: seen by A1 $\rightarrow$ A2 $\rightarrow$ A3};
  % exit arrow
  \draw[->, very thick, warnorange] (8.7,0.9) -- (1.3,0.9)
     node[midway, below=2pt, text=warnorange] {\textbf{EXIT}: seen by A3 $\rightarrow$ A2 $\rightarrow$ A1};
\end{tikzpicture}
\caption{The gate has three antennas. The order in which they detect the
ID card tells the system whether the student walked in or out.}
\end{figure}

Rooms work exactly the same way, only with \textbf{two} antennas:
antenna 1 then antenna 2 means the student \textbf{entered} the room;
antenna 2 then antenna 1 means the student \textbf{left} the room.

\begin{figure}[H]
\centering
\begin{tikzpicture}[font=\sffamily, >={Stealth[length=3mm]}]
  \fill[black!10] (0,0) rectangle (0.8,3.4);
  \fill[black!10] (7.2,0) rectangle (8.0,3.4);
  \node at (4,3.8) {\textbf{Room door (2 antennas)}};
  \foreach \x/\n in {2.8/1, 5.2/2}{
     \draw[fill=campusblue!20, draw=campusblue, thick] (\x-0.35,2.6) rectangle (\x+0.35,3.1);
     \node[text=campusblue] at (\x,2.85) {\small\textbf{A\n}};
  }
  \draw[->, very thick, okgreen] (1.0,1.7) -- (7.0,1.7)
     node[midway, above=2pt, text=okgreen] {\textbf{ENTRY}: A1 $\rightarrow$ A2};
  \draw[->, very thick, warnorange] (7.0,0.7) -- (1.0,0.7)
     node[midway, below=2pt, text=warnorange] {\textbf{EXIT}: A2 $\rightarrow$ A1};
\end{tikzpicture}
\caption{Every classroom, laboratory, library, discussion room and
auditorium door has two antennas using the same idea.}
\end{figure}

\section{What happens after a student is detected?}

\begin{enumerate}[itemsep=2pt]
  \item The moment a student walks through the \textbf{gate}, the
        parent's phone receives a notification: \emph{``Ali Khan entered
        school at 07:58''} (and the same when the student leaves).
  \item Every room the student enters during the day is silently
        recorded.
  \item Every evening (at 6:00 pm by default) the parent receives a
        \textbf{daily summary}: arrival time, departure time, rooms and
        classes attended, and any teacher remarks of the day.
  \item Every Friday the parent receives a \textbf{weekly summary} of
        the whole week.
\end{enumerate}

\section{The parts of the system}

\begin{figure}[H]
\centering
\begin{tikzpicture}[font=\sffamily\small, >={Stealth[length=3mm]},
    box/.style={draw=campusblue, fill=lightbg, rounded corners=4pt,
                minimum width=3.4cm, minimum height=1.1cm, align=center, thick}]
  \node[box] (readers) at (0,0)   {RFID readers\\ gate + all rooms};
  \node[box] (server)  at (5.6,0) {CampusTrack server\\ (.NET + MS SQL)};
  \node[box] (parent)  at (11.2,1.6)  {Parent app\\ (Android / iPhone)};
  \node[box] (student) at (11.2,0)    {Student app\\ (Android / iPhone)};
  \node[box] (portal)  at (11.2,-1.6) {Teacher \& Admin\\ web portal};
  \draw[->, thick] (readers) -- node[above]{card reads} (server);
  \draw[->, thick] (server.east) -- (parent.west);
  \draw[->, thick] (server.east) -- (student.west);
  \draw[<->, thick] (server.east) -- (portal.west);
  \node[text=black!55] at (5.6,-1.6) {\scriptsize The server stores everything and};
  \node[text=black!55] at (5.6,-1.95) {\scriptsize sends notifications to phones};
\end{tikzpicture}
\caption{How the pieces connect. The readers watch the doors, the server
keeps the records, and each person sees their own portal.}
\end{figure}

\notebox{Nobody needs to understand the technology to use CampusTrack.
Parents and students only install one app and log in. Teachers and
admins only open a website.}
